% PAPER.tex — Main file for a book/monograph project (memoir class)
%
% Compile: latexmk (shared config: ../latexmkrc); or manually
%          pdflatex PAPER; bibtex PAPER; pdflatex PAPER; pdflatex PAPER
% Engine:  pdfLaTeX
%
% SINGLE-FILE BOOK: all chapters live in this file, separated by the %# banners (navigate with outline
% folding / imenu / RefTeX; `make sep` keeps the banners tidy). A whole-book find-and-replace is one
% operation. To split a chapter out later, move it to PAPER-N-name.tex and \input it.

\documentclass[letterpaper,12pt,oneside,openany]{memoir}

% ═════════════════════════════════════════════════════════════════════════════
% ┃ ALL \usepackage, macro, and environment declarations — including the memoir ┃
% ┃ page layout (\setlrmarginsandblock, \chapterstyle, …) — live in             ┃
% ┃ PAPER.sty (loaded below), which loads the vendored shared      ┃
% ┃ packages from ../_packages/ and then this project's additions. Do NOT add       ┃
% ┃ \usepackage / \newcommand / \newtheorem / layout commands here. This        ┃
% ┃ preamble carries ONLY the document class, that one style package, and       ┃
% ┃ semantic metadata (title / author / date). See STYLE.md.                    ┃
% ═════════════════════════════════════════════════════════════════════════════
% Where the vendored shared .sty files live (one level up). This makes plain
% pdflatex work with no TEXINPUTS and no symlinks; latexmk also sets it up.
\makeatletter\def\input@path{{../_packages/}{_packages/}}\makeatother
\usepackage{PAPER}

% ─── Document Metadata ──────────────────────────────────────────────────────
\title{%
  \Huge\bfseries Title of the Monograph\\[0.5cm]
  \Large\itshape Subtitle or Series Information%
}
\author{Author One \and Author Two}
\date{\today}

% ═══════════════════════════════════════════════════════════════════════════════

%#
%# DOCUMENT: % \Huge\bfseries Title of the Monograph\\[0.5cm] \Large\itshape Subtitle or Series Information%
%#
\begin{document}

\frontmatter

% Title page
\begin{titlingpage}
  \centering
  \vspace*{4cm}
  {\Huge\bfseries Title of the Monograph\par}
  \vspace{1cm}
  {\Large\itshape Subtitle or Series Information\par}
  \vspace{3cm}
  {\large Author One \quad Author Two\par}
  \vfill
  {\large \today\par}
\end{titlingpage}

\tableofcontents*
\clearpage

\mainmatter

% ─── Content Chapters ────────────────────────────────────────────────────────
% The chapters follow, in order, in this file.
%#
%## CHAPTER: Introduction
%#
\chapter{Introduction}\label{cha:introduction}

This is the introductory chapter of the monograph.
It provides an overview of the main topics and motivates the central questions.

%#
%### SECTION: Background and Motivation
%#
\section{Background and Motivation}\label{sec:background-motivation}

The study of \ldots{} has a long history in mathematics.
We begin by recalling the fundamental definitions.

\begin{definition}\label{def:fundamental-concept}
  A \emph{widget} is a pair $(X, \tau)$ where $X$ is a set and
  $\tau \subseteq \mathcal{P}(X)$ satisfies the following axioms:
  \begin{enumerate}
  \item
    $\emptyset, X \in \tau$;
  \item
    $\tau$ is closed under arbitrary unions;
  \item
    $\tau$ is closed under finite intersections.
  \end{enumerate}
\end{definition}

%#
%### SECTION: Overview of Results
%#
\section{Overview of Results}\label{sec:overview-results}

The main result of this monograph is the following.

\begin{theorem}\label{thm:main-result}
  Every compact widget is sequentially compact.
\end{theorem}

The proof of \cref{thm:main-result} occupies \cref{cha:main-results}.

%#
%### SECTION: Organization
%#
\section{Organization}\label{sec:organization}

The monograph is organized as follows.
\Cref{cha:background} reviews the necessary prerequisites.
\Cref{cha:main-results} contains the proof of the main theorem.
\Cref{cha:conclusion} discusses open problems and future directions.
%#
%## CHAPTER: Background
%#
\chapter{Background}\label{cha:background}

This chapter collects the background material needed in subsequent chapters.

%#
%### SECTION: Notation and Conventions
%#
\section{Notation and Conventions}\label{sec:notation}

Throughout this monograph, we write $\NN$, $\ZZ$, $\QQ$, $\RR$, and $\CC$
for the natural numbers, integers, rationals, reals, and complex numbers,
respectively.

\begin{notation}\label{notn:sequences}
  A sequence in a set $X$ is a function $x \colon \NN \to X$.
  We write $(x_n)_{n \in \NN}$ or simply $(x_n)$ when the index set is clear.
\end{notation}

%#
%### SECTION: Prerequisites from Analysis
%#
\section{Prerequisites from Analysis}\label{sec:prerequisites-analysis}

We recall the basic definitions and results from analysis that will be used later.

\begin{definition}\label{def:metric-space}
  A \emph{metric space} is a pair $(X, d)$ where $X$ is a set and
  $d \colon X \times X \to [0, \infty)$ is a function satisfying:
  \begin{enumerate}
  \item
    $d(x, y) = 0$ if and only if $x = y$;
  \item
    $d(x, y) = d(y, x)$ for all $x, y \in X$;
  \item
    $d(x, z) \le d(x, y) + d(y, z)$ for all $x, y, z \in X$.
  \end{enumerate}
\end{definition}

\begin{example}\label{exa:euclidean-metric}
  The real line $\RR$ with the absolute value metric $d(x, y) = \abs{x - y}$
  is a metric space.
\end{example}

\begin{lemma}\label{lem:triangle-inequality}
  In any metric space $(X, d)$, for all $x_1, \ldots, x_n \in X$,
  \begin{equation}\label{eq:generalized-triangle}
  d(x_1, x_n) \le \sum_{i=1}^{n-1} d(x_i, x_{i+1}).
  \end{equation}
\end{lemma}

\begin{proof}
  By induction on $n$, using the triangle inequality.
\end{proof}
%#
%## CHAPTER: Main Results
%#
\chapter{Main Results}\label{cha:main-results}

In this chapter, we prove the main theorem stated in the introduction.

%#
%### SECTION: Key Lemma
%#
\section{Key Lemma}\label{sec:key-lemma}

The proof of the main theorem relies on the following key lemma.

\begin{lemma}\label{lem:key-lemma}
  Let $(X, \tau)$ be a compact widget and let $(x_n)$ be a sequence in $X$.
  Then there exists a convergent subsequence.
\end{lemma}

\begin{proof}
  \todo{Complete the proof.}
  The argument proceeds by contradiction.
  Suppose no subsequence of $(x_n)$ converges.
  Then for each $x \in X$, there exists an open neighborhood $U_x \in \tau$
  such that $U_x$ contains at most finitely many terms of the sequence.
  The collection $\{U_x : x \in X\}$ is an open cover of $X$.
  By compactness, there is a finite subcover, which leads to a contradiction.
\end{proof}

%#
%### SECTION: Proof of the Main Theorem
%#
\section{Proof of the Main Theorem}\label{sec:proof-main-theorem}

We now prove \cref{thm:main-result}.

\begin{proof}[Proof of \cref{thm:main-result}]
  This follows directly from \cref{lem:key-lemma}.
\end{proof}
%#
%## CHAPTER: Conclusion and Open Problems
%#
\chapter{Conclusion and Open Problems}\label{cha:conclusion}

We have established that every compact widget is sequentially compact
(\cref{thm:main-result}).
Several natural questions remain open.

\begin{conjecture}\label{conj:converse}
  Every sequentially compact widget is compact.
\end{conjecture}

\begin{remark}\label{rem:counterexample}
  \Cref{conj:converse} is known to hold under additional set-theoretic
  hypotheses;
  see~\cite{Kelley:1955} for details.
\end{remark}

\backmatter

% ─── Bibliography ────────────────────────────────────────────────────────────
% alpha style = [Doe25]-style labels, matching the group's other repositories.
\clearpage
\phantomsection
\addcontentsline{toc}{chapter}{\bibname}
\bibliographystyle{amsalpha}
\bibliography{../references}

\end{document}

%%% Local Variables:
%%% mode: LaTeX
%%% TeX-master: t
%%% End:
